\section{Experimental Setup}
\label{sec:setup}

We evaluate the selector family described in Section~\ref{sec:method} on two
multi-hop question-answering corpora derived from Wikipedia, each retaining
its natural hyperlink structure.
All selectors are zero-shot: no component undergoes task-specific training or
fine-tuning on the evaluation datasets.

\subsection{Datasets}
\label{sec:setup:datasets}

\paragraph{\TwoWiki.}
We use the development split of \TwoWiki~\cite{ho2020twowikimultihopqa}
as a calibration surface.
Questions require two-hop reasoning over bridge and comparison evidence
chains drawn from Wikipedia.
We draw a fixed 100-case canonical sample and evaluate under small fixed
token budgets (128 and 256 tokens).
Because most gold support sets in this dataset are compact (typically two
documents), dense retrieval is already a competitive control, making the
dataset useful for verifying method correctness and tuning hyperparameters
but insufficient as the main judgment surface.

\paragraph{\IIRC.}
The development split of \IIRC~\cite{ferguson2020iirc} serves as the primary
evaluation surface.
\IIRC questions require retrieving evidence from linked Wikipedia articles that
are \emph{not} contained in the initial context paragraph, making it a
genuinely retrieval-dependent multi-hop task.
Our full-context document store contains 61{,}304 Wikipedia articles with
their complete hyperlink structure, yielding a link-context graph
$\mathcal{G}$ with substantially higher retrieval difficulty than fragment-level
corpora.
% TODO: report |E| (total edge count) for the \IIRC graph.
We evaluate on a fixed 100-case canonical sample under ratio-controlled
selector budgets.
% TODO: confirm 100-case sample ID file after canonical \IIRC reruns.

The resulting link-context graph for \IIRC contains $|\mathcal{V}| = 61{,}304$ nodes.
% TODO: compute |E| and avg out-degree from dataset/iirc/store
Hyperlinks are extracted directly from Wikipedia article HTML; we retain all in-text hyperlinks but exclude template, infobox, and navigational links. Redirect targets are resolved to their canonical article titles. When multiple hyperlinks from the same source to the same target exist (e.g., repeated mentions), each is preserved as a separate edge with its own context triple. No anchor normalization beyond Unicode NFKC is applied.

\paragraph{Sampling protocol.}
For both datasets, we draw a fixed canonical sample of 100~cases from the development split.
Cases are selected by deterministic ordering of the dataset's native case identifiers
(no random shuffling or stratification by hop type or difficulty).
The same 100~cases are used across all selectors and budget configurations,
ensuring within-sample comparability.
% TODO: confirm exact case-ID file path after canonical IIRC reruns.

Both datasets derive from English Wikipedia and preserve natural hyperlinks
as edges in $\mathcal{G}$.
No synthetic links or knowledge-graph triples are injected; the graph
structure is exactly what appears in the source corpus.

\subsection{Metrics}
\label{sec:setup:metrics}

\paragraph{Primary metric.}
We adopt $\Fempty$ (\supportf with zero-on-empty) as defined
in Section~\ref{sec:problem:eval}.
This metric assigns $\Fempty = 0$ whenever the selected set $S$ is empty,
penalizing selectors that fail to produce any evidence rather than leaving
the score undefined.

\paragraph{Secondary metrics.}
We report support precision $\mathrm{P}(S, G)$ and support recall
$\mathrm{R}(S, G)$ to decompose selection quality, along with budget
utilization $\sum_{v \in S}\tau(v) / (\rho T)$ and budget adherence
$\mathbb{1}[\sum_{v \in S}\tau(v) \leq \rho T]$.

\paragraph{Diagnostics.}
For cost analysis we record selector runtime (wall-clock seconds),
LLM token usage (total prompt and completion tokens consumed by
LLM-based scorers), and selected node count $|S|$.
Answer-generation metrics (exact match, answer~F1) are out of scope for this
selector-first evaluation.

\subsection{Budget Configurations}
\label{sec:setup:budgets}

On \TwoWiki, we use fixed token budgets of 128 and 256 tokens for
calibration runs, matching the compact gold support sets typical of this
dataset.

On \IIRC, we use \emph{ratio-controlled} selector budgets
$\rho \in \{0.01, 0.02, 0.05, 0.10\}$, where $\rho$ is the fraction of
total corpus mass $T = \sum_{v \in \mathcal{V}} \tau(v)$ that the selector
may retain (Section~\ref{sec:problem:selection}).
We additionally include $\rho = 1.0$ as a full-corpus reference point
that subsumes all documents without selection.
Ratio-controlled budgets are the canonical setting for \IIRC because they
express selector quality relative to corpus size: at $\rho = 0.01$, the
selector must compress 61{,}304 articles down to roughly 1\% of the total
token mass, foregrounding the quality--compression tradeoff that fixed
token budgets obscure on full-document corpora.
% TODO: report the absolute token counts corresponding to each ratio
% once the canonical \IIRC store token mass is finalized.

\subsection{Selectors}
\label{sec:setup:selectors}

We compare five selectors from the family defined in
Section~\ref{sec:method:family}, plus two reference points:

\begin{itemize}[leftmargin=*,nosep]
\item \textbf{\seldense} (top-1 seed, sentence-transformer, budget
  fill): flat dense retrieval under the selector budget. This is the primary
  control and represents the strongest single-retrieval baseline available
  without graph traversal.

\item \textbf{\selmdrlight} (top-1 seed, per-hop dense expansion):
  a repo-native iterative dense baseline that approximates the
  retrieve-and-expand pattern of \MDR~\cite{xiong2021mdr} without trained
  query encoders.

\item \textbf{\selwalk} (sentence-transformer scorer,
  \stfuture profile, lookahead $L{=}2$, budget fill):
  the best non-LLM walk configuration, selected via edge-scorer ablation on
  \TwoWiki (Section~\ref{sec:results}).
  % TODO: update cross-ref if ablation table moves to analysis or appendix.

\item \textbf{\selcontroller + LLM controller}: the main
  method candidate, using an LLM-guided controller for branching and early
  stopping decisions during the walk (Section~\ref{sec:method:walk}).
  % TODO: confirm final LLM model name after 100-case \IIRC runs.

\item \textbf{\seloracle}: an oracle selector that
  returns exactly the gold support documents, establishing an upper bound on
  $\Fempty$.

\item \textbf{\selfullcorpus}: includes the entire
  corpus as the selected set ($\rho = 1.0$). This reference point achieves
  perfect recall but near-zero precision, serving as a proxy for the cost of
  passing an unselected corpus to a downstream consumer such as GraphRAG.
\end{itemize}

All selectors share the same dense seed retrieval
(\stmodel~\cite{reimers2019sentencebert}),
the same budget accounting mechanism, and the same
\budgetfill post-fill stage
(Section~\ref{sec:method:budget}), ensuring that differences in
$\Fempty$ reflect the selection strategy rather than
incidental implementation choices.

\subsection{Implementation Details}
\label{sec:setup:implementation}

Seed retrieval and embedding-based edge scoring use
\stmodel (22.7M parameters), a MiniLM variant
trained on 215M question--passage pairs~\cite{reimers2019sentencebert}.
Document embeddings are precomputed and indexed per corpus; query embeddings
are computed at selector invocation time.
% TODO: confirm whether FAISS or brute-force cosine is used at eval time.

The LLM controller in \selcontroller uses \controllerlm via the
Copilot SDK.
% TODO: update model name and report per-case token cost from 100-case runs.
We analyze sensitivity to model size in Section~\ref{sec:analysis}.
No prompt tuning or in-context examples are used.

The controller prompt presents the query, a summary of documents already selected (titles and lead sentences), and a pre-filtered set of up to 5 candidate edges with their anchor text and containing sentence. The candidate set is narrowed in two stages: first, edges with zero lexical overlap with the query are removed; second, the remaining candidates are ranked by sentence-transformer similarity and the top~5 are retained. If fewer than 2 candidates survive filtering, a rescue mechanism preserves any candidate whose anchor text has $\geq$50\% token overlap with a query entity.

The controller returns a structured JSON response containing: a ranked choice among visible candidates plus an explicit \actstop{} option, a runner-up selection for runtime branching, and natural-language reasoning. The controller does not emit per-candidate subscores directly; the decision to continue, branch (fork a scout path to the runner-up), or stop is derived from the ranked choice and runtime policy. Scout branches are limited to one per walk and are admitted only if their token cost is within 10\% of the total budget and the target document's estimated redundancy with the current selection is below~0.60.

All experiments run on a single machine with no GPU requirement for
selection (the sentence-transformer model runs on CPU).
Selector runtimes reported in the results include embedding computation,
graph traversal, and any LLM calls, measured as wall-clock time per case.
We set the walk step budget $H=2$, seed count $k=1$, score threshold $\theta=0$ (disabled), backfill pool $K=64$, and relative drop threshold $\rhofill=0.5$.
